\frfilename{mt2a3.tex}
\versiondate{25.7.07}
\copyrightdate{1995}

\def\varBbb#1{\mathchoice{\hbox{$\Bbb#1$\hskip0.02em}}
  {\hbox{$\Bbb#1$\hskip0.02em}}
  {\hbox{$\scriptstyle\Bbb#1$\hskip0.04em}}
  {\hbox{$\scriptscriptstyle\Bbb#1$\hskip0.04em}}}

\def\chaptername{Lampiran}
\def\sectionname{Topologi umum}

\newsection{2A3}

Pada berbagai bagian -- terutama \S\S245-247, %\S245 \S246 \S247
tetapi juga untuk gagasan-gagasan tertentu dalam Bab 27 -- kita perlu
mengetahui sesuatu tentang topologi yang tidak dapat dimetrikkan. Saya
menyarankan agar Anda meluangkan waktu untuk membaca sebuah buku
tentang analisis fungsional elementer yang memuat ungkapan `kekompakan lemah'
atau `kompak secara lemah' dalam indeksnya. Namun, di sini saya dapat
menyajikan konsep-konsep yang benar-benar digunakan dalam jilid ini dengan
ruang yang jauh lebih sedikit daripada yang diperlukan oleh uraian lengkap
yang ortodoks.

\leader{2A3A}{Topologi} \cmmnt{Pertama-tama kita perlu mengetahui apa
yang dimaksud dengan {\lq topologi'}.} Jika $X$ adalah sebarang himpunan,
sebuah {\bf topologi} pada $X$ adalah keluarga $\frak T$ dari himpunan bagian
$X$ sedemikian sehingga (i) $\emptyset$, $X\in\frak T$ (ii) jika $G$,
$H\in\frak T$ maka $G\cap H\in\frak T$ (iii) jika
$\Cal G\subseteq\frak T$ maka $\bigcup\Cal G\in\frak T$\cmmnt{
(bdk.\ 1A2B)}. \cmmnt{Pasangan} $(X,\frak T)$ kini merupakan sebuah
{\bf ruang topologis}. Dalam konteks ini, anggota-anggota $\frak T$ disebut
{\bf terbuka}, dan komplemennya\cmmnt{ (dalam $X$)} disebut
{\bf tertutup}\cmmnt{ (bdk.\ 1A2E-1A2F)}.

\leader{2A3B}{Fungsi kontinu (a)} Jika $(X,\frak T)$ dan $(Y,\frak
S)$ adalah ruang topologis, sebuah fungsi $\phi:X\to Y$ disebut
{\bf kontinu} jika $\phi^{-1}[G]\in\frak T$ untuk setiap
$G\in\frak S$. \cmmnt{(Menurut 2A2Ca di atas, definisi ini konsisten
dengan definisi kontinuitas $\epsilon$-$\delta$ untuk fungsi dari satu
ruang Euklides ke ruang Euklides lain. Lihat juga 2A3H di bawah.)}

\header{2A3Bb}{\bf (b)} Jika $(X,\frak T)$, $(Y,\frak S)$ dan
$(Z,\frak U)$ adalah ruang topologis, sedangkan $\phi:X\to Y$ dan
$\psi:Y\to Z$ kontinu, maka $\psi\phi:X\to Z$ kontinu.
\prooflet{\Prf\ Jika $G\in\frak U$ maka $\psi^{-1}[G]\in\frak S$, sehingga
$(\psi\phi)^{-1}[G]=\phi^{-1}[\psi^{-1}[G]]\in\frak T$.\ \Qed}

\header{2A3Bc}{\bf (c)} Jika $(X,\frak T)$ adalah ruang topologis,
sebuah fungsi $f:X\to\Bbb R$ kontinu jika dan hanya jika
$\{x:a<f(x)<b\}$ terbuka setiap kali $a<b$ dalam $\Bbb R$.
\prooflet{\Prf\ {\bf (i)} Setiap interval $\ooint{a,b}$ terbuka dalam
$\Bbb R$, sehingga jika $f$ kontinu, prabentuknya
$\{x:a<f(x)<b\}$ harus terbuka. {\bf (ii)} Misalkan
$f^{-1}[\,\ooint{a,b}\,]$ terbuka setiap kali $a<b$, dan misalkan
$H\subseteq\Bbb R$ sebarang himpunan terbuka. Menurut definisi himpunan
`terbuka' dalam $\Bbb R$\cmmnt{ (1A2A)},

$$\eqalign{H&
=\bigcup\{\ooint{y-\delta,y+\delta}:y\in\Bbb R,\,\delta>0,\,
\ooint{y-\delta,y+\delta}\subseteq H\},\cr}$$

\noindent sehingga

\Centerline{$f^{-1}[H]=
\bigcup\{f^{-1}[\,\ooint{y-\delta,y+\delta}\,]:y\in\Bbb R,\,\delta>0,\,
\ooint{y-\delta,y+\delta}\subseteq H\}$}

\noindent merupakan gabungan himpunan-himpunan terbuka dalam $X$, sehingga
terbuka.\ \Qed}

\header{2A3Bd}{\bf (d)} Jika $r\ge 1$, $(X,\frak T)$ adalah ruang
topologis, dan $\phi:X\to\BbbR^r$ adalah sebuah fungsi, maka $\phi$
kontinu jika dan hanya jika $\phi_i:X\to\Bbb R$ kontinu untuk setiap
$i\le r$, dengan $\phi(x)=(\phi_1(x),\ldots,\phi_r(x))$ untuk setiap
$x\in X$.
\prooflet{\Prf\ {\bf (i)} Misalkan $\phi$ kontinu.   Untuk $i\le
r$, $y=(\eta_1,\ldots,\eta_r)\in\BbbR^r$, tetapkan
$\pi_i(y)=\eta_i$. Maka
$|\pi_i(y)-\pi_i(z)|\le\|y-z\|$ untuk semua $y$, $z\in\BbbR^r$ sehingga
$\pi_i:\BbbR^r\to\Bbb R$
kontinu. Akibatnya $\phi_i=\pi_i\phi$ kontinu menurut (b) di atas.
{\bf (ii)} Misalkan setiap $\phi_i$ kontinu, dan
$H\subseteq\BbbR^r$ terbuka. Tetapkan

\Centerline{$\Cal G
=\{G:G\subseteq X$ terbuka, $G\subseteq\phi^{-1}[H]\}$.}

\noindent Maka $G_0=\bigcup\Cal G$ terbuka, dan
$G_0\subseteq \phi^{-1}[H]$. Namun, misalkan $x_0$ sebarang titik
dalam $\phi^{-1}[H]$. Ada $\delta>0$ sedemikian sehingga
$U(\phi(x_0),\delta)\subseteq H$, karena $H$ terbuka dan berisi
$\phi(x_0)$.   Untuk $1\le i\le r$ atur
$V_i=\{x:\phi_i(x_0)-\bover{\delta}{\sqrt{r}}
<\phi_i(x)<\phi_i(x_0)+\bover{\delta}{\sqrt{r}}\}$;
maka $V_i$ adalah gambar kebalikan dari himpunan terbuka di bawah peta kontinu
$\phi_i$, sehingga terbuka. Tetapkan $G=\bigcap_{i\le r}V_i$. Maka $G$
terbuka (menggunakan (ii) dalam definisi 2A3A), $x_0\in G$, dan
$\|\phi(x)-\phi(x_0)\|<\delta$ untuk setiap $x\in G$, sehingga
$G\subseteq\phi^{-1}[H]$, $G\in\Cal G$ dan $x_0\in G_0$. Ini
menunjukkan bahwa $\phi^{-1}[H]=G_0$ terbuka. Karena $H$ sebarang,
$\phi$ kontinu.\ \Qed}

\spheader 2A3Be Jika $(X,\frak T)$ adalah ruang topologis,
$f_1,\ldots,f_r$ adalah fungsi kontinu dari $X$ ke $\Bbb R$, dan
$h:\BbbR^r\to\Bbb R$ kontinu, maka $h(f_1,\ldots,f_r):X\to\Bbb R$
kontinu. \prooflet{\Prf\ Tetapkan
$\phi(x)=(f_1(x),\ldots,f_r(x))\in\BbbR^r$ untuk $x\in X$.   Dengan (d),
$\phi$ kontinu, jadi 2A3Bb $h(f_1,\ldots,f_r)=h\phi$ adalah
kontinu.\ \QeD} Secara khusus, $f+g$, $f\times g$ dan $f-g$ kontinu
untuk semua fungsi kontinu $f$, $g:X\to\Bbb R$.

\spheader 2A3Bf Jika $(X,\frak T)$ dan $(Y,\frak S)$ adalah ruang
topologis dan $\phi:X\to Y$ adalah fungsi kontinu, maka $\phi^{-1}[F]$
tertutup dalam $X$ untuk setiap himpunan tertutup $F\subseteq Y$.
\prooflet{(Sebab
$X\setminus\phi^{-1}[F]=\phi^{-1}[Y\setminus F]$ terbuka.)}

\leader{2A3C}{Topologi subruang}
Jika $(X,\frak T)$ adalah ruang topologis dan $D\subseteq X$,
maka $\frak T_D=\{G\cap D:G\in\frak T\}$ adalah topologi di $D$.
\prooflet{\Prf\
(i) $\emptyset=\emptyset\cap D$ dan $D=X\cap D$ milik $\frak T_D$.
(ii) Jika $G$, $H\in\frak T_D$ ada $G'$, $H'\in\frak T$ sehingga
$G=G'\cap D$, $H=H'\cap D$;  sekarang $G\cap H=G'\cap H'\cap D\in\frak T_D$.
(iii) Jika $\Cal G\subseteq\frak T_D$, tetapkan
$\Cal H=\{H:H\in\frak T,\,H\cap D\in\Cal G\}$; maka
$\bigcup\Cal G=(\bigcup\Cal H)\cap D\in\frak T_D$.\ \Qed}

$\frak T_D$ disebut {\bf topologi subruang} pada $D$, atau topologi
pada $D$ yang {\bf diinduksi} oleh $\frak T$. Jika $(Y,\frak S)$ adalah
ruang topologis lain, dan $\phi:X\to Y$ kontinu terhadap
$(\frak T,\frak S)$, maka $\phi\restr D:D\to Y$ kontinu terhadap
$(\frak T_D,\frak S)$. \prooflet{(Sebab jika $H\in\frak S$ maka

\Centerline{$(\phi\restr D)^{-1}[H]=D\cap\phi^{-1}[H]\in\frak T_D$.)}}

\leader{2A3D}{Penutupan dan interior (a)} \cmmnt{Dalam bukti 2A3Bd
saya telah menggunakan gagasan berikut.} Misalkan $(X,\frak T)$ sebarang
ruang topologis dan $A$ sebarang himpunan bagian dari $X$. Tuliskan

\Centerline{$\interior A=\bigcup\{G:G\in\frak T,\,G\subseteq A\}$.}

\noindent Maka $\interior A$ adalah\cmmnt{ himpunan terbuka, karena
merupakan gabungan himpunan-himpunan terbuka, dan tentu saja termuat dalam
$A$; dengan demikian ia harus merupakan} himpunan terbuka terbesar yang
termuat dalam $A$, dan disebut {\bf interior} dari $A$.


\header{2A3Db}{\bf (b)} Karena suatu himpunan tertutup jika dan hanya jika
komplemennya terbuka, kita mempunyai gagasan yang saling melengkapi:

\dvro{$$\eqalign{\overline{A}
&=\bigcap\{F:F\text{ tertutup},\,A\subseteq F\}
=X\setminus\interior(X\setminus A).\cr}$$}
{$$\eqalign{\overline{A}
&=\bigcap\{F:F\text{ tertutup},\,A\subseteq F\}\cr
&=X\setminus\bigcup\{X\setminus F:F\text{ tertutup},\,A\subseteq F\}\cr
&=X\setminus\bigcup\{G:G\text{ terbuka},\,A\cap G=\emptyset\}\cr
&=X\setminus\bigcup\{G:G\text{ terbuka},\,G\subseteq X\setminus A\}
=X\setminus\interior(X\setminus A).\cr}$$}

\noindent$\overline{A}$ adalah\cmmnt{ tertutup (sebagai komplemen suatu
himpunan terbuka) dan merupakan} himpunan tertutup terkecil yang memuat
$A$; himpunan ini disebut {\bf penutupan} dari $A$.
\cmmnt{(Bandingkan 2A2A.)} \cmmnt{Karena gabungan dua himpunan tertutup
adalah tertutup (bdk.\ 1A2Fc),}
$\overline{A\cup B}=\overline{A}\cup\overline{B}$ untuk semua $A$,
$B\subseteq X$.
%4{}A2G

\header{2A3Dc}{\bf (c)}\cmmnt{ Ada tak terhitung banyaknya cara untuk
memandang konsep-konsep ini; satu uraian berguna tentang penutupan suatu
himpunan adalah}

\dvro{$$\eqalign{x\in\overline A
&\iff\,\text{setiap himpunan terbuka yang memuat }x
\text{ beririsan dengan }A.\cr}$$}
{$$\eqalign{x\in\overline A
&\iff\,x\notin\interior(X\setminus A)\cr
&\iff\,\text{tidak ada himpunan terbuka yang memuat }x
\text{ dan termuat dalam }X\setminus A\cr
&\iff\,\text{setiap himpunan terbuka yang memuat }x
\text{ beririsan dengan }A.\cr}$$}

\leader{2A3E}{Topologi Hausdorff (a)}\dvro{ Sebuah}{ Konsep
`ruang topologis' dirumuskan begitu luas, dan dapat diterapkan begitu luas,
sehingga sangat banyak jenis ruang topologis yang berbeda telah dipelajari.
Untuk jilid ini kita tidak memerlukan sebagian besar kosakata (yang sangat
luas) yang telah dikembangkan untuk menggambarkan keragaman tersebut. Namun,
satu istilah yang berguna (dan salah satu konsep terpenting) adalah `ruang
Hausdorff'; sebuah} ruang topologis $X$ disebut {\bf Hausdorff} jika untuk
setiap $x$, $y\in X$ yang berbeda terdapat himpunan-himpunan terbuka saling
lepas $G$, $H\subseteq X$ sedemikian sehingga $x\in G$ dan $y\in H$.

\spheader 2A3Eb Dalam ruang Hausdorff $X$, himpunan berhingga adalah
tertutup. \prooflet{\Prf\ Jika $z\in X$, maka untuk setiap
$x\in X\setminus\{z\}$ terdapat himpunan terbuka yang memuat $x$ tetapi
tidak memuat $z$, sehingga $X\setminus\{z\}$ terbuka dan $\{z\}$ tertutup.
Jadi, suatu himpunan berhingga merupakan gabungan berhingga dari himpunan-
himpunan tertutup, sehingga tertutup.\ \Qed}

\leader{2A3F}{Pseudometrik}\cmmnt{ Banyak topologi penting (meskipun
tidak semuanya!) dapat didefinisikan oleh keluarga-keluarga pseudometrik;
akan berguna untuk memiliki sejumlah keterampilan teknis dalam hal ini.

\medskip

}{\bf (a)} Misalkan $X$ sebuah himpunan. Sebuah {\bf pseudometrik} pada
$X$ adalah fungsi $\rho:X\times X\to\coint{0,\infty}$ sedemikian sehingga

\inset{$\rho(x,z)\le\rho(x,y)+\rho(y,z)$ untuk setiap $x$, $y$, $z\in X$}

\cmmnt{\noindent (`pertidaksamaan segitiga';)}

\inset{$\rho(x,y)=\rho(y,x)$ untuk setiap $x$, $y\in X$;

$\rho(x,x)=0$ untuk setiap $x\in X$.}

\noindent Sebuah {\bf metrik} adalah pseudometrik $\rho$ yang memenuhi
syarat tambahan

\inset{jika $\rho(x,y)=0$ maka $x=y$.}

\header{2A3Fb}{\bf (b) Contoh (i)} Untuk $x$, $y\in\Bbb R$, tetapkan
$\rho(x,y)=|x-y|$\cmmnt{; maka $\rho$ adalah metrik pada $\Bbb R$}
(`metrik biasa' pada $\Bbb R$).

\medskip

\quad{\bf (ii)} Untuk $x$, $y\in\BbbR^r$, dengan $r\ge 1$, tetapkan
$\rho(x,y)=\|x-y\|$, mendefinisikan $\|z\|=\sqrt{\sum_{i=1}^r\zeta_i^2}$,
seperti biasa. Maka $\rho$ adalah\cmmnt{ sebuah metrik,} {\bf metrik
Euklides} pada $\BbbR^r$. \cmmnt{(Pertidaksamaan segitiga untuk $\rho$
berasal dari pertidaksamaan
Cauchy di 1A2C: jika $x$, $y$, $z\in\BbbR^r$, maka

\Centerline{$\rho(x,z)=\|x-z\|=\|(x-y)+(y-z)\|
\le\|x-y\|+\|y-z\|=\rho(x,y)+\rho(y,z)$.}

\noindent Sifat-sifat lain yang diperlukan dari $\rho$ bersifat elementer.
Bandingkan 2A4Bb di bawah.)}

\cmmnt{\medskip

\quad{\bf (iii)} Sebagai contoh pseudometrik yang bukan metrik,
ambil $r\ge 2$ dan definisikan $\rho:\BbbR^r\times\Bbb
R^r\to\coint{0,\infty}$ dengan menyetel $\rho(x,y)=|\xi_1-\eta_1|$ setiap kali
$x=(\xi_1,\ldots,\xi_r)$, $y=(\eta_1,\ldots,\eta_r)\in\BbbR^r$.
}%end of comment

\header{2A3Fc}{\bf (c)} Sekarang misalkan $X$ sebuah himpunan dan
$\Rho$ sebuah keluarga tak kosong dari pseudometrik pada $X$. Misalkan
$\frak T$ keluarga himpunan bagian $G$ dari $X$ sedemikian sehingga untuk
setiap $x\in G$ terdapat
$\rho_0,\ldots,\rho_n\in\Rho$ dan $\delta>0$ sedemikian rupa sehingga

\Centerline{$U(x;\rho_0,\ldots,\rho_n;\delta)
=\{y:y\in X,\,\max_{i\le n}\rho_i(y,x)<\delta\}\subseteq G$.}

\noindent Maka $\frak T$ adalah topologi pada $X$.

\prooflet{\Prf\ (Bandingkan 1A2B.) (i)
$\emptyset\in\frak T$ karena kondisinya terpenuhi secara vakum.
$X\in\frak T$ karena $U(x;\rho;1)\subseteq X$ untuk $x\in X$,
$\rho\in\Rho$.   (ii) Jika $G$, $H\in\frak T$ dan $x\in G\cap H$, ambil
$\rho_0,\ldots,\rho_m,\rho'_0,\ldots,\rho'_n\in\Rho$, $\delta$,
$\delta'>0$ sehingga $U(x;\rho_0,\ldots,\rho_m;\delta)\subseteq G$,
$U(x;\rho'_0,\ldots,\rho'_n;\delta')
\ifdim\pagewidth=390pt\penalty-1000\fi
\subseteq H$; maka


\Centerline{
$U(x;\rho_0,\ldots,\rho_m,\rho'_0,\ldots,\rho'_n;\min(\delta,\delta'))
\subseteq G\cap H$.}

\noindent Karena $x$ sebarang, $G\cap H\in\frak T$. (iii)
Jika $\Cal G\subseteq\frak T$ dan $x\in\bigcup\Cal G$, terdapat
$G\in\Cal G$ sehingga $x\in G$;  sekarang ada
$\rho_0,\ldots,\rho_n\in\Rho$ dan $\delta>0$ sedemikian sehingga


\Centerline{$U(x;\rho_0,\ldots,\rho_n;\delta)\subseteq
G\subseteq\bigcup\Cal G$.}

\noindent Karena $x$ sebarang, $\bigcup\Cal G\in\frak T$.
\Qed}

$\frak T$ adalah topologi {\bf yang didefinisikan oleh}
$\Rho$.

\header{2A3Fd}{\bf (d)} Kita dapat menggunakan konvensi untuk menangani
kasus ketika $\Rho$ adalah himpunan kosong; dalam hal ini topologi pada
$X$ yang didefinisikan oleh $\Rho$
adalah $\{\emptyset,X\}$.

\cmmnt{\spheader 2A3Fe Dalam banyak kasus penting,
$\Rho$ terarah ke atas dalam arti bahwa untuk $\rho_1$,
$\rho_2\in\Rho$ ada
$\rho\in\Rho$ sehingga $\rho_i(x,y)\le\rho(x,y)$ untuk semua $x$, $y\in X$
dan kedua nilai $i$. Dalam hal ini, tentu saja, setiap himpunan
$U(x;\rho_0,\ldots,\rho_n;\delta)$, di mana $\rho_0,\ldots,\rho_n\in\Rho$,
memuat suatu himpunan berbentuk $U(x;\rho;\delta)$, dengan $\rho\in\Rho$.
Akibatnya, misalnya, sebuah himpunan $G\subseteq X$ terbuka jika untuk setiap
$x\in G$ ada $\rho\in\Rho$, $\delta>0$ sehingga
$U(x;\rho;\delta)\subseteq G$.
}

\spheader 2A3Ff Topologi $\frak T$ disebut {\bf dapat dimetrikkan} jika
merupakan topologi yang didefinisikan oleh suatu keluarga $\Rho$ yang
terdiri dari satu metrik. Jadi topologi {\bf Euklides} pada $\BbbR^r$
dapat dimetrikkan dan didefinisikan oleh $\{\rho\}$, dengan $\rho$ metrik
dalam (b-ii) di atas.

\leader{2A3G}{Proposisi} Misalkan $X$ sebuah himpunan dengan topologi
yang didefinisikan oleh himpunan tak kosong $\Rho$ dari pseudometrik pada
$X$. Maka $U(x;\rho_0,\ldots,\rho_n;\epsilon)$ terbuka untuk setiap
$x\in X$, $\rho_0,\ldots,\rho_n\in\Rho$ dan $\epsilon>0$.

\proof{ (Bandingkan 1A2D.) Ambil
$y\in U(x;\rho_0,\ldots,\rho_n;\epsilon)$. Tetapkan

\Centerline{$\eta=\max_{i\le n}\rho_i(y,x)$,
\quad $\delta=\epsilon-\eta>0$.}

\noindent    Jika $z\in U(y;\rho_0,\ldots,\rho_n;\delta)$ maka

\Centerline{$\rho_i(z,x)\le\rho_i(z,y)+\rho_i(y,x)<\delta+\eta
=\epsilon$}

\noindent untuk setiap $i\le n$, sehingga
$U(y;\rho_0,\ldots,\rho_n;\delta)
\subseteq U(x;\rho_0,\ldots,\rho_n;\epsilon)$. Karena $y$ sebarang,
$U(x;\rho_0,\ldots,\rho_n;\epsilon)$ terbuka.
}%end of proof of 2A3G

\leader{2A3H}{}\cmmnt{ Sekarang kita mempunyai hasil yang bersesuaian
dengan 2A2Ca, yang menguraikan fungsi kontinu antara ruang topologis yang
didefinisikan oleh keluarga-keluarga pseudometrik.

\medskip

\noindent}{\bf Proposisi} Misalkan $X$ dan $Y$ himpunan; misalkan
$\Rho$ keluarga tak kosong dari pseudometrik pada $X$, dan $\Theta$
keluarga tak kosong dari pseudometrik pada $Y$; misalkan $\frak T$ dan
$\frak S$ topologi-topologi yang bersesuaian. Maka fungsi $\phi:X\to Y$
kontinu jika dan hanya jika, setiap kali $x\in X$,
$\theta\in\Theta$ dan $\epsilon>0$, ada
$\rho_0,\ldots,\rho_n\in\Rho$ dan $\delta>0$ sehingga
$\theta(\phi(y),\phi(x))\le\epsilon$ setiap kali $y\in X$ dan $\max_{i\le
n}\rho_i(y,x)\le\delta$.

\proof{{\bf (a)} Misalkan $\phi$ kontinu;  ambil $x\in X$,
$\theta\in\Theta$ dan $\epsilon>0$.   Oleh 2A3G,
$U(\phi(x);\theta;\epsilon)\in\frak S$.
Jadi $G=\phi^{-1}[U(\phi(x);\theta;\epsilon)]\in\frak T$.   Sekarang $x\in G$,
jadi ada $\rho_0,\ldots,\rho_n\in\Rho$ dan $\delta>0$ seperti itu
$U(x;\rho_0,\ldots,\rho_n;\delta)\subseteq G$.   Dalam hal ini
$\theta(\phi(y),\phi(x))\le\epsilon$ setiap kali $y\in X$ dan $\max_{i\le
n}\rho_i(y,x)\le\bover12\delta$. Karena $x$, $\theta$ dan $\epsilon$
sebarang, $\phi$ memenuhi syarat tersebut.

\medskip

{\bf (b)} Misalkan $\phi$ memenuhi syarat tersebut. Ambil
$H\in\frak S$ dan pertimbangkan $G=\phi^{-1}[H]$.   Jika $x\in G$, maka
$\phi(x)\in H$, jadi ada $\theta_0,\ldots,\theta_n\in\Theta$ dan
$\epsilon>0$ sehingga
$U(\phi(x);\theta_0,\ldots,\theta_n;\epsilon)\subseteq H$.   Untuk setiap
$i\le n$ terdapat $\rho_{i0},\ldots,\rho_{i,m_i}\in\Rho$ dan
$\delta_i>0$ sehingga $\theta_i(\phi(y),\phi(x))\le\bover12\epsilon$
setiap kali $y\in X$ dan $\max_{j\le m_i}\rho_{ij}(y,x)\le\delta_i$.   Tetapkan
$\delta=\min_{i\le n}\delta_i>0$; maka

\Centerline{$U(x;\rho_{00},\ldots,\rho_{0,m_0},\ldots,\rho_{n0},\ldots,
\rho_{n,m_n};\delta)\subseteq G$.}

\noindent Karena $x$ sebarang, $G\in\frak T$. Karena $H$ sebarang,
$\phi$ kontinu.
}%end of proof of 2A3H

\cmmnt{
\leader{2A3I}{Catatan (a)} Jika $\Rho$ terarah ke atas, syaratnya
disederhanakan menjadi: untuk setiap $x\in X$,
$\theta\in\Theta$ dan $\epsilon>0$, terdapat
$\rho\in\Rho$ dan $\delta>0$ sehingga
$\theta(\phi(y),\phi(x))\le\epsilon$ setiap kali $y\in X$ dan
$\rho(y,x)\le\delta$.




\spheader 2A3Ib Misalkan kita mempunyai sebuah himpunan $X$ dan dua
keluarga tak kosong $\Rho$, $\Theta$ dari pseudometrik pada $X$, yang
menghasilkan topologi $\frak T$ dan $\frak S$ pada $X$. Maka
$\frak S\subseteq\frak T$ jika peta identitas $\phi$ dari $X$ ke dirinya
sendiri kontinu ketika dipandang sebagai peta
dari $(X,\frak T)$ ke $(X,\frak S)$, karena ini berarti
$G=\phi^{-1}[G]$ menjadi milik $\frak T$ kapan pun $G\in\frak S$.   Menerapkan
proposisi di atas ke $\phi$, kita melihat bahwa ini terjadi jika untuk setiap
$\theta\in\Theta$,
$x\in X$ dan $\epsilon>0$ ada $\rho_0,\ldots,\rho_n\in\Rho$ dan
$\delta>0$ sehingga $\theta(y,x)\le\epsilon$ setiap kali $y\in X$ dan
$\max_{i\le n}\rho_i(y,x)\le\delta$.   Demikian pula, membalikkan peran
dari $\Rho$ dan $\Theta$, kita mendapatkan kriteria kapan $\frak
T\subseteq\frak S$, dan menggabungkan keduanya kita mendapatkan kriteria
untuk menentukan kapan $\frak T=\frak S$.
}%end of comment

\leader{2A3J}{Subruang: Proposisi} Jika $X$ adalah himpunan, $\Rho$
keluarga tak kosong dari pseudometrik pada $X$ yang mendefinisikan topologi
$\frak T$ pada $X$, dan $D\subseteq X$, maka

(a) untuk setiap $\rho\in\Rho$, pembatasan $\rho^{(D)}$ dari $\rho$
ke $D\times D$ adalah pseudometrik pada $D$;

(b) topologi yang ditentukan oleh $\Rho_D=\{\rho^{(D)}:\rho\in\Rho\}$ pada
$D$ persis dengan topologi subruang $\frak T_D$ yang dijelaskan dalam 2A3C.

\proof{(a) langsung dari definisi dalam 2A3Fa. Untuk (b), kita perlu
berpikir sejenak.

\medskip

{\bf (i)} Misalkan $G$ termasuk dalam topologi yang didefinisikan oleh
$\Rho_D$. Tetapkan


\Centerline{$\Cal H=\{H:H\in\frak T,\,H\cap D\subseteq G\}$,}

\Centerline{$H^*=\bigcup\Cal
H\in\frak T$,\quad $G^*=H^*\cap D\in\frak T_D$;}

\noindent  lalu $G^*\subseteq G$.   Di
sebaliknya, jika $x\in G$, maka ada
$\rho_0,\ldots,\rho_n\in\Rho$ dan $\delta>0$ sedemikian rupa sehingga

\Centerline{$U(x;\rho_0^{(D)},\ldots,\rho_n^{(D)};\delta)
=\{y:y\in D,\,\max_{i\le n}\rho_i^{(D)}(y,x)<\delta\}\subseteq G$.}

\noindent Pertimbangkan

\Centerline{$H=U(x;\rho_0,\ldots,\rho_n;\delta)
=\{y:y\in X,\,\max_{i\le n}\rho_i(y,x)<\delta\}\subseteq X$.}

\noindent Jelas bahwa

\Centerline{$H\cap D
=U(x;\rho_0^{(D)},\ldots,\rho_n^{(D)};\delta)\subseteq G$.}

\noindent Selain itu, $H\in\frak T$. Jadi $H\in\Cal H$ dan

\Centerline{$x\in H\cap D\subseteq H^*\cap D=G^*$.}

\noindent Dengan demikian $G=G^*\in\frak T_D$.

\medskip

{\bf (ii)} Sekarang anggaplah $G\in\frak T_D$.   Lalu ada
$H\in\frak T$ sehingga $G=H\cap D$.
Pertimbangkan peta identitas $\phi:D\to X$, yang didefinisikan dengan
$\phi(x)=x$ untuk setiap $x\in D$. Peta $\phi$ jelas memenuhi kriteria
2A3H jika kita melengkapi $D$ dengan $\Rho_D$
dan $X$ dengan $\Rho$, karena $\rho(\phi(x),\phi(y))=\rho^{(D)}(x,y)$
setiap kali $x$, $y\in D$ dan $\rho\in\Rho$; jadi $\phi$ kontinu untuk
topologi-topologi yang bersesuaian, dan $\phi^{-1}[H]$ termasuk dalam
topologi yang didefinisikan oleh $\Rho_D$. Tetapi $\phi^{-1}[H]=G$.
Jadi setiap himpunan dalam $\frak T_D$ termasuk dalam topologi yang
didefinisikan oleh $\Rho_D$, dan kedua topologi
adalah sama, seperti yang diklaim.
}%end of proof of 2A3J

\leader{2A3K}{Penutupan dan interior} Misalkan $X$ sebuah himpunan,
$\Rho$ keluarga tak kosong dari pseudometrik pada $X$, dan $\frak T$
topologi yang didefinisikan oleh $\Rho$.

\header{2A3Ka}{\bf (a)} Untuk $A\subseteq X$ dan $x\in X$,

\dvro{$$\eqalign{x\in\interior A
&\iff\,\text{terdapat }\rho_0,\ldots,\rho_n\in\Rho,\,\delta>0
\text{ sedemikian sehingga }U(x;\rho_0,\ldots,\rho_n;\delta)\subseteq A.\cr}$$}
{$$\eqalign{x\in\interior A
&\iff\,\text{terdapat himpunan terbuka yang termuat dalam }A
\text{ dan memuat }x\cr
&\iff\,\text{terdapat }\rho_0,\ldots,\rho_n\in\Rho,\,\delta>0
\text{ sedemikian sehingga }U(x;\rho_0,\ldots,\rho_n;\delta)\subseteq A.\cr}$$}

\header{2A3Kb}{\bf (b)} Untuk $A\subseteq X$ dan $x\in X$,
$x\in\overline A$ jika dan hanya jika
$U(x;\rho_0,\ldots,\rho_n;\delta)\cap A\ne\emptyset$ untuk setiap
$\rho_0,\ldots,\rho_n\in\Rho$ dan $\delta>0$.
\cmmnt{(Bandingkan 2A2B(ii), 2A3Dc.)}

\leader{2A3L}{Topologi Hausdorff}\cmmnt{ Ingat bahwa topologi
$\frak T$ adalah Hausdorff jika setiap dua titik dapat dipisahkan oleh
himpunan-himpunan terbuka (2A3E).} Topologi pada suatu himpunan $X$ yang
didefinisikan oleh keluarga tak kosong $\Rho$ dari pseudometrik adalah
Hausdorff jika dan hanya jika untuk setiap dua titik berbeda $x$, $y$
dalam $X$ terdapat $\rho\in\Rho$ sedemikian sehingga
$\rho(x,y)>0$.   \prooflet{\Prf\ {\bf (i)} Misalkan topologinya adalah
Hausdorff dan $x$, $y$ adalah titik berbeda di $X$.   Lalu ada
himpunan terbuka $G$ yang memuat $x$ tetapi tidak memuat $y$. Maka terdapat
$\rho_0,\ldots,\rho_n\in\Rho$ dan $\delta>0$ sedemikian sehingga
$U(x;\rho_0,\ldots,\rho_n;\delta)\subseteq G$, dalam hal ini
$\rho_i(y,x)\ge\delta>0$ untuk beberapa $i\le n$.   {\bf (ii)} Jika $\Rho$
memenuhi kondisi, dan $x$, $y$ adalah titik berbeda dari $X$, ambil
$\rho\in\Rho$ sehingga $\rho(x,y)>0$, dan atur
$\delta=\bover12\rho(x,y)$. Maka $U(x;\rho;\delta)$ dan
$U(y;\rho;\delta)$ saling lepas (sebab jika $z\in X$, maka

\Centerline{$\rho(z,x)+\rho(z,y)\ge\rho(x,y)=2\delta$,}

\noindent sehingga setidaknya salah satu dari $\rho(z,x)$, $\rho(z,y)$
lebih besar dari atau sama dengan $\delta$), dan keduanya merupakan
himpunan terbuka yang masing-masing memuat $x$, $y$. Karena $x$ dan $y$
sebarang, topologinya Hausdorff.
\Qed}

Secara khusus, setiap topologi yang dapat dimetrikkan adalah Hausdorff.

\leader{2A3M}{Konvergensi barisan (a)} Jika $(X,\frak T)$ adalah
sebarang ruang topologis, dan $\sequencen{x_n}$ adalah barisan dalam $X$,
kita mengatakan bahwa $\sequencen{x_n}$ {\bf konvergen} ke $x\in X$,
atau bahwa $x$ adalah {\bf limit} dari $\sequencen{x_n}$, atau
$\sequencen{x_n}\to x$, jika untuk setiap himpunan terbuka $G$ yang
memuat $x$ terdapat $n_0\in\Bbb N$ sehingga
$x_n\in G$ untuk setiap $n\ge n_0$.

\header{2A3Mb}{\bf (b) Peringatan} Dalam ruang topologis umum,
sebuah barisan mungkin mempunyai lebih dari satu limit\cmmnt{, sehingga
kita tidak dapat menuliskan $x=\lim_{n\to\infty}x_n$} dengan aman. Namun,
dalam ruang Hausdorff hal ini tidak terjadi. \prooflet{\Prf\ Jika
$\frak T$ Hausdorff, dan $x$,
$y$ adalah titik berbeda dari $X$, terdapat himpunan terbuka yang terpisah $G$, $H$
sehingga $x\in G$ dan $y\in H$.   Jika sekarang $\sequencen{x_n}$ menyatu ke
$x$, terdapat $n_0$ sehingga $x_n\in G$ untuk setiap $n\ge n_0$, sehingga
$x_n\notin H$ untuk setiap $n\ge n_0$, dan $\sequencen{x_n}$ tidak dapat
konvergen ke $y$.\ \Qed} Secara khusus, suatu barisan dalam ruang yang
dapat dimetrikkan mempunyai paling banyak satu limit.

\header{2A3Mc}{\bf (c)} Misalkan $X$ sebuah himpunan dan $\Rho$ keluarga
tak kosong dari pseudometrik pada $X$, yang menghasilkan topologi
$\frak T$; misalkan $\sequencen{x_n}$ barisan dalam $X$ dan $x\in X$.
Maka $\sequencen{x_n}$ konvergen ke $x$ jika dan hanya jika
$\lim_{n\to\infty}\rho(x_n,x)=0$
untuk setiap $\rho\in\Rho$.   \prooflet{\Prf\ {\bf (i)} Misalkan
$\sequencen{x_n}\to x$ dan $\rho\in\Rho$. Maka untuk setiap
$\epsilon>0$, himpunan $G=U(x;\rho;\epsilon)$ terbuka dan memuat $x$,
sehingga terdapat $n_0$ dengan $x_n\in G$ untuk setiap $n\ge n_0$,
yakni $\rho(x_n,x)<\epsilon$ untuk setiap $n\ge n_0$. Karena $\epsilon$
sebarang, $\lim_{n\to\infty}\rho(x_n,x)=0$. {\bf (ii)} Jika syaratnya
terpenuhi, ambil himpunan terbuka $G$ yang memuat $x$. Maka
ada $\rho_0,\ldots,\rho_k\in\Rho$ dan $\delta>0$ sedemikian sehingga
$U(x;\rho_0,\ldots,\rho_k;\delta)\subseteq G$. Untuk setiap $i\le k$
terdapat $n_i\in\Bbb N$ sedemikian sehingga $\rho_i(x_n,x)<\delta$ untuk
setiap $n\ge n_i$. Tetapkan $n^*=\max(n_0,\ldots,n_k)$; maka
$x_n\in U(x;\rho_0,\ldots,\rho_k;\delta)\subseteq G$ untuk setiap
$n\ge n^*$. Karena $G$ sebarang, $\sequencen{x_n}\to x$.\ \Qed}

\spheader 2A3Md Misalkan $(X,\rho)$ ruang metrik, $A$ himpunan bagian
dari $X$, dan $x\in X$. Maka $x\in\overline{A}$ jika dan hanya jika
terdapat barisan dalam $A$ yang konvergen ke $x$.
\prooflet{\Prf {\bf (i)} Jika $x\in\overline{A}$,
maka untuk setiap $n\in\Bbb N$ kita dapat memilih satu titik
$x_n\in A\cap U(x;\rho;2^{-n})$
(2A3Kb);  sekarang $\sequencen{x_n}\to x$.   {\bf (ii)} Jika $\sequencen{x_n}$
adalah barisan dalam $A$ yang konvergen ke $x$, maka untuk setiap himpunan terbuka $G$
yang berisi $x$ terdapat $n$ sehingga $x_n\in G$, sehingga $A\cap
G\ne\emptyset$;  oleh 2A3Dc, $x\in\overline{A}$.\ \Qed}

\leader{2A3N}{Kekompakan} \cmmnt{ Konsep berikutnya yang kita perlukan
adalah gagasan `kekompakan' dalam ruang topologis umum.

\header{2A3Na}}{\bf (a)} Jika $(X,\frak T)$ sebarang ruang topologis,
sebuah himpunan bagian $K$ dari $X$ disebut {\bf kompak} jika, setiap kali
$\Cal G$ adalah keluarga dalam $\frak T$ yang menutupi $K$, terdapat
keluarga berhingga $\Cal G_0\subseteq\Cal G$ yang menutupi $K$.
\cmmnt{(Bdk.\ 2A2D. Sebuah {\bf peringatan}: banyak penulis mencadangkan
istilah `kompak' untuk ruang Hausdorff.)} Sebuah himpunan $A\subseteq X$
disebut {\bf relatif kompak} dalam $X$ jika terdapat himpunan bagian
kompak dari $X$ yang memuat $A$. \cmmnt{({\bf Peringatan!} Dalam ruang
non-Hausdorff, ini tidak sama dengan mengatakan bahwa
$\overline{A}$ kompak.)}

\header{2A3Nb}{\bf (b)}\cmmnt{ Sebagaimana dalam 2A2E-2A2G %2A2E 2A2F 2A2G
(dan bukti-buktinya sama dalam kasus umum), kita memperoleh hasil-hasil
berikut.

\medskip

\quad}{\bf (i)} Jika $K$ kompak dan $E$ tertutup, maka $K\cap E$ adalah
kompak.

\medskip

\quad {\bf (ii)} Jika $K\subseteq X$ kompak dan $\phi:K\to Y$ adalah
kontinu,
dengan $(Y,\frak S)$ ruang topologis lain, maka $\phi[K]$ adalah himpunan
bagian kompak dari $Y$.

\medskip

\quad{\bf (iii)} Jika $K\subseteq X$ kompak dan $\phi:K\to\Bbb R$ adalah
kontinu, maka $\phi$ terbatas dan mencapai nilai maksimum serta minimumnya.

\leader{2A3O}{Titik gugus (a)} Jika $(X,\frak T)$ adalah ruang
topologis dan $\sequencen{x_n}$ adalah barisan dalam $X$, maka sebuah
{\bf titik gugus} dari $\sequencen{x_n}$ adalah titik $x\in X$ sedemikian
sehingga, setiap kali $G$ merupakan himpunan terbuka yang memuat $x$ dan
$n\in\Bbb N$, terdapat $k\ge n$ dengan $x_k\in G$.

\header{2A3Ob}{\bf (b)} Jika $(X,\frak T)$ adalah ruang topologis dan
$A\subseteq X$ relatif kompak, maka setiap barisan $\sequencen{x_n}$
dalam $A$ mempunyai titik gugus dalam $X$.
\prooflet{\Prf\
Misalkan $K$ himpunan bagian kompak dari $X$ yang memuat $A$. Tetapkan

\Centerline{$\Cal G=\{G:G\in\frak T,\,\{n:x_n\in G\}\text{
berhingga}\}$.}

\noindent\Quer\ Jika $\Cal G$ mencakup $K$, maka ada
$\Cal G_0\subseteq\Cal G$ berhingga yang mencakup $K$.   Sekarang

\Centerline{$\Bbb N=\{n:x_n\in A\}=\{n:x_n\in\bigcup\Cal G_0\}
=\bigcup_{G\in\Cal G_0}\{n:x_n\in G\}$}

\noindent merupakan gabungan berhingga dari himpunan-himpunan berhingga,
suatu kemustahilan. \Bang\ Jadi $\Cal G$ tidak menutupi $K$. Ambil
sebarang $x\in K\setminus\bigcup\Cal G$. Jika $G\in\frak T$, $x\in G$
dan $n\in\Bbb N$, maka $G\notin\Cal G$, sehingga $\{k:x_k\in G\}$ tak
berhingga dan terdapat $k\ge n$ sedemikian sehingga $x_k\in G$. Jadi $x$
adalah titik gugus dari $\sequencen{x_n}$, sebagaimana diperlukan.\ \Qed}

\cmmnt{
\leader{2A3P}{Filter} Dalam $\BbbR^r$, dan lebih umum dalam semua ruang
yang dapat dimetrikkan, gagasan-gagasan topologis dapat dibahas secara
efektif dengan barisan konvergen. (Memang, hal ini kadang-kadang memerlukan
suatu bentuk lemah dari aksioma pilihan untuk memilih sebuah barisan; tetapi
karena teori ukuran tanpa pilihan semacam itu berubah sama sekali -- lihat
Bab 56 dalam Jilid 5 -- tidak ada gunanya mempermasalahkannya di sini.)
Namun, untuk ruang topologis secara umum, barisan sama sekali tidak memadai,
karena alasan-alasan sangat menarik yang tidak akan saya uraikan. Sebagai
gantinya kita perlu menggunakan `jaring' atau `filter'. Yang terakhir
memerlukan sedikit usaha lebih besar pada awalnya, tetapi kemudian (menurut
pendapat saya) jauh lebih mudah digunakan; karena itu saya menguraikan
metode ini sekarang.
}%end of comment

\leader{2A3Q}{Filter konvergen (a)} Misalkan $(X,\frak T)$ ruang
topologis, $\Cal F$ filter pada $X$\cmmnt{ (lihat 2A1I)}, dan $x$ titik
dalam $X$. Kita mengatakan bahwa $\Cal F$ {\bf konvergen} ke $x$, atau
bahwa $x$ adalah {\bf limit} dari $\Cal F$, dan menuliskan
$\Cal F\to x$, jika setiap himpunan terbuka yang memuat $x$ merupakan
anggota $\Cal F$.

\header{2A3Qb}{\bf (b)} Misalkan $(X,\frak T)$ dan $(Y,\frak S)$ ruang
topologis, $\phi:X\to Y$ fungsi kontinu, $x\in X$, dan $\Cal F$ filter
pada $X$ yang konvergen ke $x$. Maka $\phi[[\Cal F]]$\cmmnt{ (sebagaimana
didefinisikan dalam 2A1Ib)} konvergen ke $\phi(x)$\prooflet{ (sebab
$\phi^{-1}[G]$ adalah himpunan terbuka yang memuat $x$ setiap kali $G$
merupakan himpunan terbuka yang memuat $\phi(x)$)}.

\leader{2A3R}{}\cmmnt{ Sekarang kita mempunyai karakterisasi kekompakan
berikut.

\medskip

\noindent}{\bf Teorema} Misalkan $X$ ruang topologis dan $K$ himpunan
bagian dari $X$. Maka $K$ kompak jika dan hanya jika setiap ultrafilter
pada $X$ yang memuat $K$ mempunyai limit dalam $K$.

\proof{{\bf (a)} Misalkan $K$ kompak dan $\Cal F$
adalah ultrafilter pada $X$ yang berisi $K$.   Tetapkan

\Centerline{$\Cal G=\{G:G\subseteq X$ terbuka,
$X\setminus G\in\Cal F\}$.}

\noindent Maka gabungan setiap dua anggota $\Cal G$ merupakan anggota
$\Cal G$, sehingga gabungan sejumlah berhingga anggota $\Cal G$ juga
merupakan anggota $\Cal G$; selain itu, tidak ada anggota $\Cal G$ yang
dapat memuat $K$, karena $X\setminus K\notin\Cal F$. Karena $K$ kompak,
$\Cal G$ tidak dapat menutupi $K$. Misalkan $x$ sebarang titik dalam
$K\setminus\bigcup\Cal G$. Jika $G$ sebarang himpunan terbuka yang memuat
$x$, maka $G\notin\Cal G$, sehingga $X\setminus G\notin\Cal F$; tetapi
ini berarti $G$ harus merupakan anggota $\Cal F$, karena $\Cal F$ adalah
ultrafilter. Karena $G$ sebarang, $\Cal F\to x$. Jadi setiap ultrafilter
pada $X$ yang memuat $K$ mempunyai limit dalam $K$.

\medskip

{\bf (b)} Sekarang anggap bahwa setiap ultrafilter pada $X$ yang memuat
$K$ mempunyai limit dalam $K$. Misalkan $\Cal G$ penutup $K$ oleh
himpunan-himpunan terbuka dalam $X$. \Quer\ Andaikan, jika mungkin,
$\Cal G$ tidak mempunyai subpenutup berhingga. Tetapkan

\Centerline{$\Cal F=\{F:$ terdapat keluarga berhingga sedemikian sehingga
$\Cal G_0\subseteq\Cal G,\,F\cup\bigcup\Cal G_0\supseteq K\}$.}

\noindent Maka $\Cal F$ adalah filter pada $X$. \Prf\ (i)
$X\cup\bigcup\emptyset\supseteq K$, sehingga $X\in\Cal F$.


\Centerline{$\emptyset\cup\bigcup\Cal G_0=\bigcup\Cal G_0\not\supseteq
K$}

\noindent untuk setiap keluarga berhingga $\Cal G_0\subseteq\Cal G$,
menurut hipotesis, sehingga
$\emptyset\notin\Cal F$.
(ii) Jika $E$, $F\in\Cal F$ terdapat himpunan berhingga $\Cal G_1$,
$\Cal G_2\subseteq\Cal G$ sehingga $E\cup\bigcup\Cal G_1$ dan
$F\cup\bigcup\Cal G_2$ keduanya menyertakan $K$;  sekarang
$(E\cap F)\cup\bigcup(\Cal G_1\cup\Cal G_2)\supseteq K$ jadi
$E\cap F\in\Cal F$.   (iii) Jika
$X\supseteq E\supseteq F\in\Cal F$ maka terdapat
$\Cal G_0\subseteq\Cal G$ berhingga sehingga
$F\cup\bigcup\Cal G_0\supseteq K$; sekarang
$E\cup\bigcup\Cal G_0\supseteq K$ dan $E\in\Cal F$.\ \Qed

Menurut Teorema Ultrafilter (2A1O), terdapat ultrafilter $\Cal F^*$ pada
$X$ yang memuat $\Cal F$. Tentu saja $K$ sendiri merupakan anggota
$\Cal F$, sehingga $K\in\Cal F^*$. Menurut hipotesis, $\Cal F^*$
mempunyai limit $x\in K$. Tetapi sekarang terdapat himpunan $G\in\Cal G$
yang memuat $x$, dan
$(X\setminus G)\cup G\supseteq K$, jadi
$X\setminus G\in\Cal F\subseteq\Cal F^*$;  yang berarti
bahwa $G$ tidak dapat merupakan anggota $\Cal F^*$, dan $x$ tidak dapat
menjadi limit dari $\Cal F^*$. \Bang

Jadi $\Cal G$ mempunyai subpenutup berhingga. Karena $\Cal G$ sebarang,
$K$ harus kompak.
}%end of proof of 2A3R

\cmmnt{
\medskip

\noindent{\bf Catatan} Perhatikan bahwa bagian (b) dari pembuktian
teorema ini sangat bergantung pada Teorema Ultrafilter
dan oleh karena itu pada aksioma pilihan.
}

\leader{2A3S}{Perhitungan lebih lanjut dengan filter (a)}
S\cmmnt{ecara umum, s}ebuah filter mungkin mempunyai lebih dari satu
limit; tetapi dalam ruang Hausdorff hal ini tidak terjadi.
\prooflet{\Prf\ (Bandingkan 2A3Mb.) Jika
$(X,\frak T)$ adalah Hausdorff, dan $x$,
$y$ adalah titik berbeda dalam $X$, terdapat himpunan-himpunan terbuka
saling lepas $G$, $H$
sehingga $x\in G$ dan $y\in H$.   Jika sekarang filter $\Cal F$ di $X$
menyatu ke
$x$, maka $G\in\Cal F$, sehingga $H\notin\Cal F$ dan $\Cal F$ tidak
konvergen ke
$y$.\ \Qed}

Karena itu, kita dapat menuliskan $x=\lim\Cal F$ dengan aman ketika
$\Cal F\to x$ dalam ruang Hausdorff.

\spheader 2A3Sb Sekarang misalkan $X$ sebuah himpunan, $\Cal F$ filter
pada $X$, $(Y,\frak S)$ ruang Hausdorff, $D\in\Cal F$, dan
$\phi:D\to Y$ sebuah fungsi. Kita menuliskan
$\lim_{x\to\Cal F}\phi(x)$ untuk $\lim\phi[[\Cal F]]$ jika limit ini
terdefinisi dalam $Y$; yakni $\lim_{x\to\Cal F}\phi(x)=y$ jika dan hanya
jika $\phi^{-1}[H]\in\Cal F$ untuk setiap himpunan terbuka $H$ yang
memuat $y$.

\dvAnew{2014}Jika $Z$ adalah himpunan lain, $\Cal G$ filter pada $Z$,
dan $\psi:Z\to X$ sedemikian sehingga $\Cal F=\psi[[\Cal G]]$, maka
komposisi $\phi\psi$ terdefinisi pada $\psi^{-1}[D]\in\Cal G$; jika
salah satu dari limit $\lim_{x\to\Cal F}\phi(x)$ dan
$\lim_{z\to\Cal G}\phi\psi(z)$ terdefinisi dalam $Y$, maka yang lain juga
terdefinisi, dan keduanya sama. \prooflet{\Prf\ Misalkan $y\in Y$ dan
$\Cal U$ keluarga himpunan bagian terbuka dari $Y$ yang memuat $y$. Maka

$$\eqalign{\lim_{x\to\Cal F}\phi(x)=y
&\iff\phi^{-1}[G]\in\Cal F\text{ untuk setiap }G\in\Cal U
\cr&\iff\psi^{-1}[\phi^{-1}[G]]\in\Cal G\text{ untuk setiap }G\in\Cal U
\cr&\iff(\phi\psi)^{-1}[G]\in\Cal G\text{ untuk setiap }G\in\Cal U
\iff\lim_{z\to\Cal G}\phi\psi(z)=y. \text{ \Qed}\cr}$$}

Dalam kasus khusus $Y=\Bbb R$, $\lim_{x\to\Cal F}\phi(x)=a$ jika dan
hanya jika
$\{x:|\phi(x)-a|\le\epsilon\}\in\Cal F$ untuk setiap
$\epsilon>0$\prooflet{ (karena setiap himpunan terbuka yang memuat $a$
memuat suatu himpunan berbentuk $[a-\epsilon,a+\epsilon]$, yang pada
gilirannya memuat himpunan terbuka
 $\ooint{a-\epsilon,a+\epsilon}$)}.

\spheader 2A3Sc Misalkan $X$ dan $Y$ himpunan, $\Cal F$ filter pada
$X$, $\Theta$ keluarga tak kosong dari pseudometrik pada $Y$ yang
mendefinisikan topologi $\frak S$ pada $Y$, dan $\phi:X\to Y$ sebuah
fungsi. Maka filter citra $\phi[[\Cal F]]$ konvergen ke $y\in Y$ jika dan
hanya jika $\lim_{x\to\Cal F}\theta(\phi(x),y)=0$ dalam
$\Bbb R$
untuk setiap $\theta\in\Theta$.   \prooflet{\Prf\ {\bf (i)} Misalkan
$\phi[[\Cal F]]\to
y$.   Untuk setiap $\theta\in\Theta$ dan $\epsilon>0$,
$U(y;\theta;\epsilon)=\{z:\theta(z,y)<\epsilon\}$ adalah himpunan terbuka
yang berisi $y$ (2A3G), jadi milik $\phi[[\Cal F]]$, dan gambar
kebalikannya $\{x:0\le\theta(\phi(x),y)<\epsilon\}$ milik $\Cal F$.
Karena $\epsilon$ bersifat arbitrer, $\lim_{x\to\Cal F}\theta(\phi(x),y)=0$.
Karena $\theta$ bersifat sewenang-wenang, $\phi$ memenuhi ketentuan tersebut.   {\bf (ii)}
Sekarang
misalkan $\lim_{x\to\Cal F}\theta(\phi(x),y)=0$ untuk setiap
$\theta\in\Theta$.   Misalkan $G$ menjadi himpunan terbuka apa pun di $Y$ yang berisi $y$.
Lalu ada $\theta_0,\ldots,\theta_n\in\Theta$ dan $\epsilon>0$ seperti
yang

\Centerline{$U(y;\theta_0,\ldots,\theta_n;\epsilon)
=\bigcap_{i\le n}U(y;\theta_i;\epsilon)\subseteq G$.}

\noindent Untuk setiap $i\le n$,

\Centerline{$\phi^{-1}[U(y;\theta_i;\epsilon)]
=\{x:\theta_i(\phi(x),y)<\epsilon\}$}

\noindent merupakan anggota $\Cal F$; karena $\Cal F$ tertutup terhadap
irisan berhingga, demikian pula
$\phi^{-1}[U(y;\theta_0,\ldots,\theta_n;\epsilon)]$ dan superhimpunannya
$\phi^{-1}[G]$. Jadi $G\in\phi[[\Cal F]]$. Karena $G$ sebarang,
$\phi[[\Cal F]]\to y$.\ \Qed}

%2A3S

\spheader 2A3Sd Secara khusus,\cmmnt{ dengan mengambil $X=Y$ dan $\phi$
sebagai peta identitas,} jika $X$ mempunyai topologi $\frak T$ yang
didefinisikan oleh keluarga tak kosong $\Rho$ dari pseudometrik, maka
filter $\Cal F$ pada $X$ konvergen ke $x\in X$ jika dan hanya jika
$\lim_{y\to\Cal F}\rho(y,x)=0$ untuk setiap $\rho\in\Rho$.

\spheader 2A3Se{\bf (i)} Jika $X$ adalah himpunan apa pun, $\Cal F$ adalah ultrafilter
pada $X$, $(Y,\frak S)$ adalah ruang Hausdorff, dan
$h:X\to Y$ adalah fungsi sedemikian sehingga $h[F]$ relatif kompak dalam
$Y$ untuk suatu $F\in\Cal F$, maka $\lim_{x\to\Cal F}h(x)$ terdefinisi
dalam $Y$. \prooflet{\Prf\ Misalkan $K\subseteq Y$ himpunan kompak yang
memuat $h[F]$. Maka $K\in h[[\Cal F]]$, yang merupakan ultrafilter
(2A1N), sehingga $h[[\Cal F]]$ mempunyai limit dalam $Y$ (2A3R), yaitu
$\lim_{x\to\Cal F}h(x)$.\ \Qed}

\medskip

\quad{\bf (ii)} Jika $X$ adalah himpunan apa pun, $\Cal F$ adalah ultrafilter pada $X$,
dan $h:X\to\Bbb R$ adalah fungsi sehingga $h[F]$ dibatasi dalam $\Bbb R$
untuk suatu himpunan $F\in\Cal F$, maka $\lim_{x\to\Cal F}h(x)$ ada dalam
$\Bbb R$.
\prooflet{\Prf\ $\overline{h[F]}$ bersifat tertutup dan dibatasi, oleh karena itu
kompak (2A2F), sehingga $h[F]$ relatif kompak dan kita dapat menggunakan
(i).\ \Qed}

\spheader 2A3Sf \cmmnt{ Konsep $\limsup$, $\liminf$ dapat diterapkan
ke filter.}  Misalkan $\Cal F$ adalah filter pada himpunan $X$,
dan $f:X\to[-\infty,\infty]$ adalah fungsi apa pun.   Kemudian

\dvro{$$\eqalign{\limsup_{x\to\Cal F}f(x)
&=\inf_{F\in\Cal F}\sup_{x\in F}f(x)
\in[-\infty,\infty],\cr}$$}
{$$\eqalign{\limsup_{x\to\Cal F}f(x)
&=\inf\{u:u\in[-\infty,\infty],\,\{x:f(x)\le u\}\in\Cal F\}\cr
&=\inf_{F\in\Cal F}\sup_{x\in F}f(x)
\in[-\infty,\infty],\cr}$$}

\dvro{$$\eqalign{\liminf_{x\to\Cal F}f(x)
&=\sup_{F\in\Cal F}\inf_{x\in F}f(x).\cr}$$}
{$$\eqalign{\liminf_{x\to\Cal F}f(x)
&=\sup\{u:u\in[-\infty,\infty],\,\{x:f(x)\ge u\}\in\Cal F\}\cr
&=\sup_{F\in\Cal F}\inf_{x\in F}f(x).\cr}$$}

\noindent\dvro{Untuk}{Mudah dilihat bahwa, untuk} sebarang dua fungsi $f$,
$g:X\to\Bbb R$,

\Centerline{$\lim_{x\to\Cal F}f(x)=a$\quad jika dan hanya jika\quad
$a=\limsup_{x\to\Cal F}f(x)=\liminf_{x\to\Cal F}f(x)$,}

\noindent dan

\Centerline{$\limsup_{x\to\Cal F}(f(x)+g(x))
\le\limsup_{x\to\Cal F}f(x)+\limsup_{x\to\Cal F}g(x)$,}

\Centerline{$\liminf_{x\to\Cal F}(f(x)+g(x))
\ge\liminf_{x\to\Cal F}f(x)+\liminf_{x\to\Cal F}g(x)$,}

\Centerline{$\liminf_{x\to\Cal F}(-f(x))=-\limsup_{x\to\Cal F}f(x)$,
\quad$\limsup_{x\to\Cal F}(-f(x))=-\liminf_{x\to\Cal F}f(x)$,}

\Centerline{$\liminf_{x\to\Cal F}cf(x)=c\liminf_{x\to\Cal F}f(x)$,
\quad$\limsup_{x\to\Cal F}cf(x)=c\limsup_{x\to\Cal F}f(x)$}

\noindent setiap kali sisi kanan ditentukan di
$[-\infty,\infty]$ dan $c\ge 0$.   Jadi jika $a=\lim_{x\to\Cal F}f(x)$ dan
$b=\lim_{x\to\Cal F}g(x)$ ada dalam $\Bbb R$, maka
$\lim_{x\to\Cal F}(f(x)+g(x))$ ada dan sama dengan $a+b$, sedangkan
$\lim_{x\to\Cal F}cf(x)$ ada dan
sama dengan $c\lim_{x\to\Cal F}f(x)$ untuk setiap $c\in\Bbb R$.

\dvro{Jika}{Kita juga melihat bahwa jika} $f:X\to\Bbb R$ sedemikian sehingga

\Centerline{untuk setiap $\epsilon>0$ terdapat $F\in\Cal F$ sedemikian sehingga
$\sup_{x\in F}f(x)\le\epsilon+\inf_{x\in F}f(x)$,}

\noindent maka
\cmmnt{$\limsup_{x\to\Cal F}f(x)\le\epsilon+\liminf_{x\to\Cal F}f(x)$
untuk setiap $\epsilon>0$, sehingga} $\lim_{x\to\Cal F}f(x)$ didefinisikan dalam
$[-\infty,\infty]$.

\spheader 2A3Sg \cmmnt{Perhatikan bahwa limit-limit standar dalam analisis
real dapat dinyatakan dalam bentuk yang diuraikan di sini. Misalnya,}
$\lim_{n\to\infty}$, $\limsup_{n\to\infty}$, $\liminf_{n\to\infty}$
berhubungan dengan $\lim_{n\to\CalFr}$, $\limsup_{n\to\CalFr}$,
$\liminf_{n\to\CalFr}$ di mana $\CalFr$ adalah filter {\bf Fr\'echet}
pada $\Bbb N$, yakni filter $\{\Bbb N\setminus A:A\subseteq\Bbb N$
berhingga$\}$ dari himpunan-himpunan bagian $\Bbb N$ yang berkomplemen
berhingga. Demikian pula,
$\lim_{\delta\downarrow a}$, $\limsup_{\delta\downarrow a}$,
$\liminf_{\delta\downarrow a}$ berhubungan dengan $\lim_{\delta\to\Cal F}$,
$\limsup_{\delta\to\Cal F}$, $\liminf_{\delta\to\Cal F}$ di mana

\Centerline{$\Cal F=\{A:A\subseteq\Bbb R,\,\exists\,h>0$ sedemikian sehingga
$\ocint{a,a+h}\subseteq A\}$.}
%2A3S

\leader{2A3T}{Topologi produk}\cmmnt{ Kita memerlukan beberapa catatan singkat
mengenai topologi pada ruang hasil kali.

\medskip

} {\bf (a)} Misalkan $(X,\frak T)$ dan $(Y,\frak S)$ adalah ruang topologis.
Misalkan $\frak U$ adalah himpunan semua himpunan bagian $U$ dari $X\times Y$
sedemikian sehingga untuk setiap $(x,y)\in U$ terdapat $G\in\frak T$ dan
$H\in\frak S$ dengan
$(x,y)\in G\times H\subseteq U$.   Maka $\frak U$ merupakan topologi pada
$X\times Y$.   \prooflet{\Prf\ (i) $\emptyset\in\frak U$ karena syarat
keanggotaan dalam $\frak U$ terpenuhi secara vakum.   $X\times Y\in\frak U$
karena $X\in\frak T$, $Y\in\frak S$ dan $(x,y)\in X\times Y\subseteq
X\times Y$ untuk setiap $(x,y)\in X\times Y$.   (ii) Jika $U$, $V\in\frak U$
dan $(x,y)\in U\cap V$, maka terdapat $G$, $G'\in\frak T$, $H$,
$H'\in\frak S$ sehingga

\Centerline{$(x,y)\in G\times H\subseteq U$,
\quad $(x,y)\in G'\times H'\subseteq V$;}

\noindent sekarang $G\cap G'\in \frak T$, $H\cap H'\in\frak S$ dan

\Centerline{$(x,y)\in(G\cap G')\times(H\cap H')\subseteq U\cap V$.}

\noindent Karena $(x,y)$ sebarang, $U\cap V\in\frak U$.   (iii) Jika
$\Cal U\subseteq\frak U$ dan $(x,y)\in\bigcup\Cal U$, maka terdapat
$U\in\Cal U$ sehingga $(x,y)\in U$;  sekarang ada $G\in\frak T$,
$H\in\frak S$ sedemikian sehingga $(x,y)\in G\times H\subseteq
U\subseteq\bigcup\Cal U$.   Karena $(x,y)$ bersifat arbitrer, $\bigcup\Cal
U\in\frak U$.\ \Qed}

$\frak U$ disebut {\bf topologi produk} pada $X\times Y$.


\spheader 2A3Tb Misalkan, dalam (a), $\frak T$ dan $\frak S$ ditentukan oleh
keluarga tak kosong $\Rho$ dan $\Theta$ dari pseudometrik\cmmnt{ dengan cara
2A3F}.   Maka $\frak U$ ditentukan oleh keluarga
$\Upsilon=\{\tilde\rho:\rho\in\Rho\}\cup\{\bar\theta:\theta\in\Theta\}$
pseudometrik pada $X\times Y$, dengan

\Centerline{$\tilde\rho((x,y),(x',y'))=\rho(x,x')$,
\quad$\bar\theta((x,y),(x',y'))=\theta(y,y')$}

\noindent setiap kali $x$, $x'\in X$, $y$, $y'\in Y$, $\rho\in\Rho$ dan
$\theta\in\Theta$.

\prooflet{\Prf\ {\bf (i)} Tentu saja kita harus memeriksa bahwa setiap
$\tilde\rho$, $\bar\theta$ adalah pseudometrik pada $X\times Y$.

\quad{\bf (ii)} Jika $U\in\frak U$ dan $(x,y)\in U$, maka ada
$G\in\frak T$, $H\in\frak S$ sehingga $(x,y)\in G\times H\subseteq U$.
Ada $\rho_0,\ldots,\rho_m\in\Rho$,
$\theta_0,\ldots,\theta_n\in\Theta$, $\delta$, $\delta'>0$ sedemikian sehingga
(dalam bahasa 2A3Fc) $U(x;\rho_0,\ldots,\rho_m;\delta)\subseteq G$,
$U(y;\theta_0,\ldots,\theta_n;\delta')\subseteq H$.   Sekarang

\Centerline{$U((x,y);\tilde\rho_0,\ldots,\tilde\rho_m,
\bar\theta_0,\ldots,\bar\theta_n;\min(\delta,\delta'))\subseteq U$.}

\noindent Karena $(x,y)$ sebarang, $U$ terbuka dalam topologi yang dihasilkan
oleh $\Upsilon$.

\quad{\bf (iii)} Jika $U\subseteq X\times Y$ terbuka dalam topologi yang
ditentukan oleh $\Upsilon$, ambil sebarang $(x,y)\in U$.   Maka ada
$\upsilon_0,\ldots,\upsilon_k\in\Upsilon$ dan $\delta>0$ sehingga
$U((x,y);\upsilon_0,\ldots,\upsilon_k;\delta)\subseteq U$.   Ambil
$\rho_0,\ldots,\rho_m\in\Rho$ dan $\theta_0,\ldots,\theta_n\in\Theta$
sehingga $\{\upsilon_0,\ldots,\upsilon_k\}
\subseteq\{\tilde\rho_0,\ldots,\tilde\rho_m,
\bar\theta_0,\ldots,\bar\theta_n\}$;    lalu
$G=U(x;\rho_0,\ldots,\rho_m;\delta)\in\frak T$ (2A3G),
$H=U(y;\theta_0,\ldots,\theta_n;\delta)\in\frak S$, dan

\Centerline{$(x,y)\in G\times H
=U((x,y);\tilde\rho_0,\ldots,\tilde\rho_m,
\bar\theta_0,\ldots,\bar\theta_n;\delta)
\subseteq U((x,y);\upsilon_0,\ldots,\upsilon_k;\delta)\subseteq U$.}

\noindent Karena $(x,y)$ sebarang, $U\in\frak U$.   Ini melengkapi bukti bahwa
$\frak U$ adalah topologi yang ditentukan oleh $\Upsilon$.\ \Qed}

\spheader 2A3Tc Secara khusus, topologi produk pada
$\BbbR^r\times\BbbR^s$ adalah topologi Euklides jika kita mengidentifikasi
$\BbbR^r\times\BbbR^s$ dengan $\BbbR^{r+s}$.  \prooflet{\Prf\ Topologi produk
ditentukan oleh dua pseudometrik $\upsilon_1$ dan $\upsilon_2$; untuk $x$,
$x'\in \BbbR^r$ dan $y$, $y'\in\BbbR^s$, kita tulis

\Centerline{$\upsilon_1((x,y),(x',y'))=\|x-x'\|$,
\quad$\upsilon_2((x,y),(x',y'))=\|y-y'\|$}

\noindent (2A3F(b-ii)).   Demikian pula, topologi Euklides pada
$\BbbR^r\times\BbbR^s\cong\BbbR^{r+s}$ ditentukan oleh metrik $\rho$,
di mana

\Centerline{$\rho((x,y),(x',y'))=\|(x,y)-(x',y')\|
=\sqrt{\|x-x'\|^2+\|y-y'\|^2}$.}

\noindent Sekarang jika $(x,y)\in\BbbR^r\times\BbbR^s$ dan $\epsilon>0$, maka

\Centerline{$U((x,y);\rho;\epsilon)
\subseteq U((x,y);\upsilon_j;\epsilon)$}

\noindent untuk kedua nilai $j$, sedangkan

\Centerline{$U((x,y);\upsilon_1,\upsilon_2;\Bover{\epsilon}{\sqrt2})
\subseteq U((x,y);\rho;\epsilon)$.}

\noindent Jadi, seperti dikemukakan dalam 2A3Ib, masing-masing topologi
termuat dalam yang lain, sehingga keduanya sama.\ \Qed}

\leader{2A3U}{Himpunan padat (a)} Jika $X$ adalah ruang topologis, himpunan
$D\subseteq X$ disebut {\bf padat} dalam $X$ jika $\overline{D}=X$\cmmnt{,
yaitu jika setiap himpunan terbuka tak kosong beririsan dengan $D$}.   Secara
lebih umum, jika $D\subseteq A\subseteq X$, maka $D$ padat dalam $A$ jika himpunan itu
padat untuk topologi subruang pada $A$\cmmnt{ (2A3C), yaitu jika
$A\subseteq\overline{D}$}.   %for 253

\spheader 2A3Ub Jika $\frak T$ ditentukan oleh keluarga tak kosong $\Rho$ dari
pseudometrik pada $X$, maka $D\subseteq X$ padat jika dan hanya jika
\discrcenter{468pt}
{$U(x;\rho_0,\ldots,\rho_n;\delta)\cap D\ne\emptyset$ }setiap kali
$x\in X$, $\rho_0,\ldots,\rho_n\in\Rho$ dan $\delta>0$.

\spheader 2A3Uc Jika $(X,\frak T)$ dan $(Y,\frak S)$ adalah ruang topologis,
dengan $Y$ Hausdorff (khususnya, jika $(X,\rho)$
dan $(Y,\theta)$ adalah ruang metrik),
dan $f$, $g:X\to Y$ adalah fungsi kontinu yang sama nilainya pada suatu
himpunan bagian padat $D$ dari $X$, maka $f=g$.
\prooflet{\Prf\Quer\ Misalkan, jika memungkinkan, terdapat $x\in X$
sehingga $f(x)\ne g(x)$.   Maka ada himpunan terbuka $G$, $H\subseteq Y$
sedemikian sehingga $f(x)\in G$, $g(x)\in H$ dan $G\cap H=\emptyset$.
Sekarang $f^{-1}[G]\cap g^{-1}[H]$ adalah himpunan terbuka yang memuat $x$
dan karena itu tak kosong, tetapi tidak beririsan dengan $D$; jadi
$x\notin\overline{D}$ dan $D$ tidak padat.\ \Bang\QeD}

\spheader 2A3Ud Ruang topologis disebut {\bf separabel} jika memiliki himpunan
bagian padat yang terhitung.   \cmmnt{Misalnya, $\BbbR^r$ separabel untuk
setiap $r\ge 1$, karena $\varBbb{Q}^r$ padat.}

\discrpage
